% Part: turing-machines
% Chapter: machines-computations
% Section: unary-numbers

\documentclass[../../../include/open-logic-section]{subfiles}

\begin{document}

\olfileid{tur}{mac}{una}
\olsection{संख्यांचे एकचिन्ही निरूपण}

\begin{explain}
ट्यूरिंग यंत्रे त्यांच्या फितीवर लिहिलेल्या चिन्हांच्या क्रमिकांवर
कार्य करतात. यंत्र वापरत असलेल्या वर्णमालेनुसार या क्रमिका विविध
आदानांचे आणि प्रदानांचे निरूपण करू शकतात. विशेष रस आहे तो
\emph{अंकगणितीय} फलने, म्हणजे नैसर्गिक संख्यांची फलने, संगणित करणाऱ्या
ट्यूरिंग यंत्रांत. धन पूर्णांकांचे निरूपण करण्याचा एक सोपा मार्ग म्हणजे
एकाच चिन्हाच्या~$\TMstroke$ क्रमिकांनी त्यांचे सांकेतीकरण करणे.
$n \in \Nat$ असेल, तर $n = 0$ असताना $\TMstroke^n$ ही रिकामी
क्रमिका मानू; अन्यथा ती नेमक्या $n$ $\TMstroke$ चिन्हांची क्रमिका
असेल.
\end{explain}

\begin{defn}[संगणन]
ट्यूरिंग यंत्र~$M$ हे $f\colon \Nat^k \to \Nat$ फलन \emph{संगणित
करते}, असे तेव्हा आणि केवळ तेव्हाच म्हणतो, जेव्हा पुढील स्वरूपाचे
प्रत्येक आदान घेतल्यावर
\[
\TMstroke^{n_1} \TMblank \TMstroke^{n_2} \TMblank \dots \TMblank \TMstroke^{n_k}
\]
$M$~थांबते आणि त्याचे प्रदान $\TMstroke^{f(n_1, \dots, n_k)}$ असते.
\end{defn}

\begin{prob}
ट्यूरिंग यंत्र~$M$ हे $f\colon \Nat^k \to \Nat^m$ फलन
कधी संगणित करते याची व्याख्या द्या.
\end{prob}

\begin{ex}\ollabel{ex:adder}
\emph{बेरीज:}
$f(n,m) = n + m$ हे फलन संगणित करणारे यंत्र बनवू.
यासाठी सुरुवातीला फितीवर अनुक्रमे $n$ आणि~$m$ लांबीचे
$\TMstroke$ चिन्हांचे दोन खंड असले पाहिजेत आणि शेवटी
$n+m$~$\TMstroke$ चिन्हांचा एकच खंड राहिल्यावर यंत्र थांबले
पाहिजे. दोन $\TMstroke$ आदान-खंडांमध्ये~$\TMblank$ असते. म्हणून एक पद्धत अशी:
$\TMblank$ असलेल्या घरात एक रेघ लिहा आणि शेवटचे~$\TMstroke$
पुसून टाका.
\begin{figure}\[
\begin{tikzpicture}[->,>=stealth',shorten >=1pt,auto,node distance=2.8cm,
                    semithick]
  \tikzstyle{every state}=[fill=none,draw=black,text=black]

  \node[initial,state]         (A)              {$q_0$};
  \node[state]         (B) [right of=A] {$q_1$};
  \node[state]         (C) [right of=B] {$q_2$};

  \path (A) edge node {\TMtrans{\TMblank}{\TMstroke}{\TMstay}} (B)
           edge [loop above] node {\TMtrans{\TMstroke}{\TMstroke}{\TMright}} (B)
        (B) edge [loop above] node {\TMtrans{\TMstroke}{\TMstroke}{\TMright}} (B)
            edge node {\TMtrans{\TMblank}{\TMblank}{\TMleft}} (C)
        (C) edge [loop above] node {\TMtrans{\TMstroke}{\TMblank}{\TMstay}} (C);
\end{tikzpicture}
\]\caption{$f(x,y) = x+y$ संगणित करणारे यंत्र}
\ollabel{fig:adder}
\end{figure}
\end{ex}

\begin{prob}
$\tuple{3,2}$ हे आदान असताना \olref[tur][mac][una]{ex:adder}
मधील यंत्राच्या स्थितिवर्णनांचा क्रम पाहा. यंत्राने $0+0$ संगणित
केल्यास काय घडते?
\end{prob}

\begin{explain}
\olref[rep]{ex:doubler} मध्ये आपण एका ट्यूरिंग यंत्राचे उदाहरण दिले:
ते~$\TMstroke$ चिन्हांची क्रमिका आदान म्हणून घेते आणि फितीवर
दुप्पट~$\TMstroke$ चिन्हांची क्रमिका ठेवून थांबते. हे दुप्पट करणारे
यंत्र आहे. पण दुप्पट झालेल्या $\TMstroke$ खंडाच्या डावीकडेही
प्रदानात $\TMblank$ चिन्हे येतात. म्हणून, वाटू शकते तसे, ते
प्रत्यक्षात $f(x) = 2x$ हे फलन संगणित करत नाही. ही अडचण
दूर करण्याचे दोन मार्ग पाहू.
\end{explain}

\begin{ex}
\olref{fig:doubler-disc} मधील यंत्र $f(x) =
2x$ हे फलन संगणित करते. \olref[rep]{ex:doubler} मधील यंत्र
प्रत्येक आदानातील~$\TMstroke$ पुसून फितीच्या अगदी उजवीकडे दोन
$\TMstroke$ चिन्हे लिहिते. त्याऐवजी हे यंत्र आदानातील प्रत्येक
$\TMstroke$ साठी उजवीकडे एकच~$\TMstroke$ वाढवते. आदान कुठे
संपते हे लक्षात ठेवण्यासाठी ते आदान आणि वाढवलेल्या रेघा यांच्यामध्ये
एक~$\TMblank$ ठेवते; अगदी शेवटी त्या जागी~$\TMstroke$ लिहिते.
तसेच आदानात सध्या कुठे आहोत हे ``लक्षात ठेवण्यासाठी'' आदान-खंडातील
एक~$\TMstroke$ तात्पुरते~$\TMblank$ ने बदलते.
  \begin{figure}
\[
\begin{tikzpicture}[->,>=stealth',shorten >=1pt,auto,node distance=2.8cm,
                    semithick]
  \tikzstyle{every state}=[fill=none,draw=black,text=black]

  \node[initial,state] (0)              {$q_0$};
  \node[state]         (1) [above of=0] {$q_1$};
  \node[state]         (2) [above of=1] {$q_2$};
  \node[state]         (3) [right of=2] {$q_3$};
  \node[state]         (4) [below of=3] {$q_4$};
  \node[state]         (5) [below of=4] {$q_5$};
  \node[state]         (6) [above right of=4] {$q_6$};
  \node[state]         (7) [below of=6] {$q_7$};
  \node[state]         (8) [below of=7] {$q_8$};

  \path (0) edge node {\TMtrans{\TMstroke}{\TMblank}{\TMright}} (1)
%    (0) edge node {\TMtrans{\TMblank}{\TMblank}{\TMstay}} (h)
    (1) edge [loop left] node {\TMtrans{\TMstroke}{\TMstroke}{\TMright}} (1)
      edge node {\TMtrans{\TMblank}{\TMblank}{\TMright}} (2)
    (2) edge [loop above] node {\TMtrans{\TMstroke}{\TMstroke}{\TMright}} (2)
        edge node {\TMtrans{\TMblank}{\TMstroke}{\TMleft}} (3)
    (3) edge [loop above] node {\TMtrans{\TMstroke}{\TMstroke}{\TMleft}} (3)
        edge node[left] {\TMtrans{\TMblank}{\TMblank}{\TMleft}} (4)
    (4) edge        node[left] {\TMtrans{\TMstroke}{\TMstroke}{\TMleft}} (5)
    (5) edge [loop below] node {\TMtrans{\TMstroke}{\TMstroke}{\TMleft}} (5)
        edge              node {\TMtrans{\TMblank}{\TMstroke}{\TMright}} (0)
    (4) edge      node[sloped] {\TMtrans{\TMblank}{\TMstroke}{\TMright}} (6)
    (6) edge              node {\TMtrans{\TMblank}{\TMstroke}{\TMright}} (7)
    (7) edge [loop right] node {\TMtrans{\TMstroke}{\TMstroke}{\TMright}} (7)
        edge              node {\TMtrans{\TMblank}{\TMblank}{\TMleft}} (8)
    (8) edge [loop right] node {\TMtrans{\TMstroke}{\TMblank}{\TMstay}} (8);
    \end{tikzpicture}
\]
\caption{$f(x) = 2x$ संगणित करणारे यंत्र}
\ollabel{fig:doubler-disc}
\end{figure}
\end{ex}

\begin{ex}\ollabel{ex:mover}
$f(x) = 2x$ संगणित करण्याचा दुसरा मार्ग म्हणजे मूळ दुप्पट
करणारे यंत्र तसेच ठेवून, शेवटी दुप्पट झालेला रेघांचा खंड फितीच्या
अगदी डावीकडे नेणाऱ्या अवस्था आणि सूचना जोडणे. \olref{fig:mover}
मधील यंत्र नेमका हा शेवटचा भाग करते: सुरुवातीला फितीवर
प्रथम~$\TMblank$ चिन्हांचा आणि नंतर~$\TMstroke$ चिन्हांचा खंड
असतो (आणि शीर्ष~$\TMblank$ खंडातील कोणत्याही घरावर असते);
मग ते $\TMstroke$ चिन्हे एकेक करून पुसते आणि फितीच्या सुरुवातीला
लिहिते. काम कधी संपले हे ओळखता यावे म्हणून ते प्रथम
$\TMstroke$ खंडाचा शेवट~$\TMendtape$ चिन्हाने दाखवते;
शेवटी हे चिन्ह पुसले जाते. अवस्था~$q_6$ पासून मोजायला
सुरुवात केली आहे, जेणेकरून त्या मूळ दुप्पट करणाऱ्या यंत्रात जोडता
येतील. त्यासाठी एक अतिरिक्त सूचना
$\delta(q_0,
\TMblank) = \tuple{q_6,\TMblank,\TMstay}$, म्हणजे~$q_0$ पासून
~$q_6$ कडे~$\TMtrans{\TMblank}{\TMblank}{\TMstay}$ हे लेबल असलेला
बाण, एवढीच हवी. (येथे एक सूक्ष्म अडचण आहे: परिणामी यंत्र
आदान~$x=0$ साठी चालत नाही. ही अडचण सरावासाठी ठेवू.)
\begin{figure}
  \[
  \begin{tikzpicture}[->,>=stealth',shorten >=1pt,auto,node distance=2.8cm,
                      semithick]
    \tikzstyle{every state}=[fill=none,draw=black,text=black]
    \node[initial,state] (6)              {$q_6$};
    \node[state]         (7) [right of=6] {$q_7$};
    \node[state]         (8) [right of=7] {$q_8$};
    \node[state]         (9) [below of=8] {$q_9$};
    \node[state]         (10) [left of=9] {$q_{10}$};
    \node[state]         (11) [left of=10]  {$q_{11}$};
    \node[state]         (12) [below of=11]  {$q_{12}$};
    \node[state]         (13) [right of=12] {$q_{13}$};
    \node[state]         (14) [right of=13] {$q_{14}$};
    \path
    (6)  edge node {\TMtrans{\TMblank}{\TMblank}{\TMright}} (7)
    (7)  edge [loop above] node {\TMtrans{\TMblank}{\TMblank}{\TMright}} (7)
         edge node {\TMtrans{\TMstroke}{\TMstroke}{\TMright}} (8)
    (8)  edge [loop above] node {\TMtrans{\TMstroke}{\TMstroke}{\TMright}} (8)
         edge node {\TMtrans{\TMblank}{\TMendtape}{\TMleft}} (9)
    (9)  edge [loop right] node {\TMtrans{\TMstroke}{\TMstroke}{\TMleft}} (9)
         edge node {\TMtrans{\TMblank}{\TMblank}{\TMright}} (10)
    (10) edge node[above] {\TMtrans{\TMstroke}{\TMblank}{\TMleft}} (11)
    (11) edge [loop above] node {\TMtrans{\TMblank}{\TMblank}{\TMleft}} (11)
         edge node[left] {\begin{tabular}{@{}l@{}}
          \TMtrans{\TMendtape}{\TMendtape}{\TMright}\\
          \TMtrans{\TMstroke}{\TMstroke}{\TMright}
         \end{tabular}} (12)
    (12) edge node[] {\TMtrans{\TMblank}{\TMstroke}{\TMright}} (13)
    (13) edge [loop below] node {\TMtrans{\TMblank}{\TMblank}{\TMright}} (13)
         edge node[sloped] {\TMtrans{\TMstroke}{\TMblank}{\TMleft}} (11)
    (13) edge node {\TMtrans{\TMendtape}{\TMblank}{\TMstay}} (14);
  \end{tikzpicture}
  \]
  \caption{$\TMstroke$ चिन्हांचा खंड डावीकडे नेणे}
  \ollabel{fig:mover}
\end{figure}
\end{ex}

\begin{prob}
\olref[tur][mac][una]{ex:mover} मध्ये आपण
\olref[tur][mac][una]{fig:doubler-disc} मधील दुप्पट करणारे यंत्र
आणि \olref[tur][mac][una]{fig:mover} मधील खंड हलवणारे यंत्र
एकत्र करून बनणाऱ्या यंत्राचे वर्णन केले. हे संयुक्त यंत्र
आदान~$x=0$, म्हणजे रिकामी फीत, घेऊन सुरू केल्यास काय घडते?
या वेळी यंत्राने प्रदान~$2x=0$ देऊन थांबावे यासाठी त्यात
काय बदल कराल? (एक अवस्था आणि एक संक्रमण जोडून हे करता
आले पाहिजे.)
\end{prob}

\begin{prob}
\emph{वजाबाकी:} $n$ आणि~$m$ लांबीच्या रेघांच्या दोन अरिक्त
क्रमिका आदान म्हणून दिल्या असताना, जेथे $n >
m$ आहे, $f(n,m) = n - m$ हे फलन संगणित करणारे ट्यूरिंग यंत्र
रचा.
\end{prob}

\begin{prob}
\emph{समता:} पुढील फलन संगणित करणारे ट्यूरिंग यंत्र रचा:
\[
\fn{equality}(n,m) = 
\begin{cases}
  \text{1} & \text{$n = m$ असेल तर} \\
  \text{0} & \text{$n \neq m$ असेल तर}
\end{cases}
\]
येथे~$n$ आणि~$m \in \PosInt$ आहेत.
\end{prob}

\begin{prob}
$x$ आणि $y$ हे धन पूर्णांक फितीवर $\TMblank$ ने विभक्त केलेल्या
$\TMstroke$ चिन्हांच्या क्रमिकांनी निरूपित केले असताना
$\min(x,y)$ हे फलन संगणित करणारे ट्यूरिंग यंत्र रचा.
यंत्राच्या वर्णमालेत अतिरिक्त चिन्हे वापरू शकता.

$\min$ हे फलन आपल्या आदानांतील लहान मूल्य निवडते; म्हणून
$\min(3,5)=3$, $\min(20,16)=16$, $\min(4,4)=4$, इत्यादी.
\end{prob}

\begin{defn}
ट्यूरिंग यंत्र~$M$ हे $f\colon \Nat^k
  \pto \Nat$ हे आंशिक फलन तेव्हा आणि केवळ तेव्हाच संगणित
करते, जेव्हा
  \begin{enumerate}
    \item $f(n_1, \dots, n_k) = m$ असेल तर $M$ हे
    $\TMstroke^{n_1}\concat\TMblank\concat
    \dots \concat\TMblank\concat\TMstroke^{n_k}$ आदानावर
    थांबते आणि त्याचे प्रदान $\TMstroke^{m}$ असते.
    \item $f(n_1, \dots, n_k)$ अपरिभाषित असेल तर $M$
    मुळीच थांबत नाही, किंवा थांबले तरी त्याचे प्रदान
    $\TMstroke$ चिन्हांचा एकच खंड नसते.
  \end{enumerate}
\end{defn}

\end{document}
